\documentclass[11pt,a4paper]{article}

\usepackage{fontspec}
\usepackage[english]{babel}
\defaultfontfeatures{Ligatures=TeX}
\setmainfont{Latin Modern Roman}
\setsansfont{Latin Modern Sans}
\setmonofont{Latin Modern Mono}

\usepackage[unicode,colorlinks=true,linkcolor=black,citecolor=blue,urlcolor=blue]
  {hyperref}
\usepackage{amsfonts,amsmath}
\usepackage{booktabs}
\usepackage{csquotes}
\usepackage{etoolbox}
\usepackage{geometry}
\usepackage{graphicx}
\usepackage{setspace}
\usepackage{xcolor}

\geometry{
  tmargin=3.0cm,
  bmargin=3.0cm,
  lmargin=2.8cm,
  rmargin=2.8cm,
  footskip=24pt
}

\setcounter{secnumdepth}{3}
\setlength{\parindent}{0pt}
\setlength{\parskip}{0.7\baselineskip}

\newcommand{\todo}[1]{{\color{red}#1}}

\def\UNIT#1#2{\ifstrempty{#2}{}{\ifstrequal{#2}{1}{\mathrm{#1}}{\mathrm{#1}^{#2}}}}
\def\units#1#2#3{\ifstrempty{#1#2#3}{$[-]$}{$[ \UNIT{kg}{#1}\UNIT{m}{#2}\UNIT{s}{#3} ]$}}
\def\unitss#1#2#3{\ifstrempty{#1#2#3}{$-$}{$ \UNIT{kg}{#1}\UNIT{m}{#2}\UNIT{s}{#3} $}}
\def\munit#1{\,\mathrm{#1}}
\def\bunit#1{\left[\mathrm{#1}\right]}

\def\abs#1{\lvert#1\rvert}
\def\argdot{{\hspace{0.18em}\cdot\hspace{0.18em}}}
\def\d{\mathrm{d}}
\def\D{{\tn D}}
\def\div{\operatorname{div}}
\def\grad{\nabla}
\def\jmp#1{[#1]}
\def\prtl{\partial}
\def\R{\mathbb R}
\def\Real{\mathbb R}
\def\sc#1#2{\left(#1,#2\right)}
\def\tn#1{{\mathbb{#1}}}
\def\vc#1{\mathbf{\boldsymbol{#1}}}
\def\E{{\mathsf E}}
\def\avg#1{\langle#1\rangle}
\def\var#1{\langle\!\langle#1\rangle\!\rangle}
\def\Var{\mathop{\rm Var}}
\def\Cov{\mathop{\rm Cov}}

\title{Template Title for a Scientific Paper or Technical Report}
\author{Author Name \\ Institution or Project}
\date{\today}

\begin{document}
\begin{onehalfspacing}

\maketitle

\begin{abstract}
Replace this abstract with a concise statement of the problem, the method, the
main result, and the scope of the document.
\end{abstract}

\tableofcontents
\clearpage

\section{Introduction}
\label{sec:introduction}

This repository provides a minimal starting point for scientific papers and
technical reports. Replace the template text with project-specific content and
keep the structure only where it supports the final document.

State the document goal early. For a paper, frame the research question or
technical claim. For a report, state the practical purpose, scope, and intended
use of the results.

\section{Methods or Model}
\label{sec:methods}

Describe the method, data, or model with enough detail for review and reuse.
Define notation before relying on it repeatedly. For example, let
$u \colon \Omega \to \R$ denote a scalar field on a domain $\Omega \subset \R^d$.
Then a generic balance equation may be written as
\[
  -\div\left(\tn D \grad u\right) = f \qquad \text{in } \Omega.
\]

The template keeps a small set of reusable macros for vectors, tensors,
operators, and units. Use project macros consistently instead of reintroducing
raw notation locally.

\section{Results}
\label{sec:results}

Summarize the key result in a form the reader can inspect quickly. Tables and
figures should support a specific claim rather than act as a general archive.

\begin{table}[htb]
  \centering
  \caption{Example table with a technical caption that can stand on its own.}
  \label{tab:example-results}
  \begin{tabular}{lcc}
    \toprule
    Case & Value & Unit \\
    \midrule
    Reference & 1.00 & $-$ \\
    Variant A & 1.12 & $-$ \\
    Variant B & 0.94 & $-$ \\
    \bottomrule
  \end{tabular}
\end{table}

Reference tables, figures, and equations through labels such as
Tab.~\ref{tab:example-results} and Section~\ref{sec:methods}.

\section{Discussion}
\label{sec:discussion}

Interpret the results without overstating them. Make assumptions, limitations,
and sensitivity to inputs explicit where they affect the conclusions.

\section{Conclusion}
\label{sec:conclusion}

Keep the conclusion traceable to the preceding sections. Summarize what was
shown, what remains uncertain, and what a reader should take forward.

\end{onehalfspacing}
\end{document}
